\section{Estrutura de Intervalos (Counting Process)}

Para modelar a recorrência e os riscos concorrentes, o dataset deve ser reestruturado para o formato de \textbf{Processo de Contagem (t\_start, t\_stop)}. Cada empresa é representada por múltiplos intervalos que indicam a permanência ou transição entre estados.

\begin{table}[h]
\centering
\caption{Exemplo de Estrutura de Dados Multi-State}
\begin{tabular}{@{}ccccccc@{}}
\toprule
\textbf{NIF} & \textbf{t\_start} & \textbf{t\_stop} & \textbf{From\_State} & \textbf{To\_State} & \textbf{Status} & \textbf{Episode (k)} \\ \midrule
123 & 0 & 3 & 0 & 1 & 1 & 1 \\
123 & 3 & 5 & 1 & 0 & 1 & 1 \\
123 & 5 & 8 & 0 & 2 & 1 & 2 \\ \bottomrule
\end{tabular}
\end{table}

\section{Algoritmo de Definição de Estados}

A lógica de transição segue os critérios de negócio definidos:
\begin{enumerate}
    \item \textbf{Eventos Terminais:} Se o \textit{Status} legal indica falência (C3), transição imediata para o Estado 2 (Absorvente).
    \item \textbf{Distress Transiente:} Se ocorrer insolvência técnica (C1: Equity < 0) ou incapacidade de cobertura (C2: EBITDA < Juros), transição para o Estado 1.
    \item \textbf{Recuperação:} Se uma empresa em Estado 1 deixa de verificar C1 e C2, regressa ao Estado 0.
\end{enumerate}

\section{Prevenção de Data Leakage: Lagging e Censura}

\begin{itemize}
    \item \textbf{Lagging Temporal:} Todas as variáveis independentes ($X$) devem ser defasadas em 1 ano ($t-1$) em relação ao estado alvo ($y_t$), refletindo a disponibilidade real da informação contabilística.
    \item \textbf{Gestão de Censura:} Empresas que deixam de reportar dados sem indicação de falência são marcadas como censuradas à direita no seu último registo ativo.
\end{itemize}
